\section{Relation to channel metrology and measurement information}
\label{sec:comparison84}

\subsection{Support spans, channel estimation and finite pairs}
For the measurement Kraus operators
$L_{j,a}=|j\rangle\langle a|\sqrt{E_j}$, the real Hermitian
Kraus-product span is exactly $\mathcal W_E$.  For the orbit
$L_{j,a}(t)=L_{j,a}e^{-itG}$ the effective generator is $G$.
Zhou--Jiang \cite[Theorems~1--3 and Section~V]{ZJ84} establish the
Kraus-span criterion for the linear or quadratic asymptotic channel
Fisher-information rate and give error-correction attainability.
We do not claim a new metrological dichotomy or a new correction
principle.  Theorems~\ref{thm:supportkernel84} and
\ref{thm:orbittrichotomy84} identify its complete support-kernel
form and give finite-angle trace-distance inequalities with an
explicit $4mt^2\|G\|_{\mathrm{op}}^2$ remainder.
Theorem~\ref{thm:curveclassification85} adds a normalized first-order
realization for general two-sided $C^2$ effect curves; its
$\sqrt N|t|+O(Nt^2)$ bound handles rank changes without a smooth
exact Kraus hypothesis.  Theorem~\ref{thm:controlstability85} separately
budgets certified control errors.  Neither constitutes a new claim
to the underlying metrological exponents.

Sieniawski--Demkowicz-Dobrza\'nski
\cite[Section~IV, equations~(16)--(23)]{SD76} integrate channel
Fisher-information bounds to bound adaptive discrimination and
compare metrological scaling with perfect finite discrimination.
The same horizontal method is a principal antecedent of
\eqref{eq:frontupper84}.  Our support-kernel identity is stated in
measurement coordinates and our lower \eqref{eq:finiteanglelower84}
comes from an actual corrected finite-angle channel.  It does not
claim the exact optimal tester or perfect discrimination.


\subsection{Strategy classes and finite-use comparisons}
Our independent inputs are products of \emph{complete} probe--reference
systems.  Public mixtures preserve the bound; a correlated reference
shared between groups does not satisfy this definition.  Block-$b$
preparations interpolate between those inputs and unrestricted
entangled parallel inputs.  Classical feedback permits earlier
labels to change later preparations and is not a parallel strategy;
fully adaptive control additionally retains quantum memory.  The
independent-block theorem is not a classification of every feedback
model.  Theorem~\ref{thm:resetwidth88} treats the separately defined
classical-feedback reset class with arbitrary receiver memory; it
still excludes coherent import into a fresh block.

Harrow--Hassidim--Leung--Watrous \cite[Sections~3--4]{HHLW87}
construct channels perfectly distinguishable by two adaptive calls
but not by any finite entangled parallel strategy.  Those are
perfect-discrimination statements about a fixed pair of general
channels.  Corollary~\ref{cor:parallel87} concerns comparison orders
in shrinking fixed-tangent neighborhoods of classical-output
measurements; it asserts neither equality of distances nor perfect
discrimination.  There is no contradiction with that separation.

Salek--Hayashi--Winter \cite[Theorem~9 and Remark~10]{SHW87}
compare asymptotic Chernoff and Hoeffding rates for classical
feed-forward with no quantum memory at the input against nonadaptive
product strategies.  Their restrictions and fixed-pair error
exponents differ from the present hard call cap and finite trace
separation.  In particular their product result is not an upper
bound on our entangled parallel inputs.

Li--Hirche--Tomamichel \cite[Theorems~3.1--3.3]{LHT87}
characterize sequential error-exponent regions under expected and
high-probability stopping constraints in their specified adaptive,
block and nonadaptive models, under finite max-relative-entropy
hypotheses.  Our budget instead bounds the number of calls on every
record, and our parameter $s$ varies jointly with that budget.
Their general instrument discussion is not invoked as a theorem of
unrestricted quantum-memory optimality here.

Bounds in terms of sums of squared entangled group sizes, and GHZ
attainability of metrological enhancements, are standard
\cite{Hyllus78,Toth78}.  Theorem~\ref{thm:widthlaw87} does not
claim these mechanisms as new.  It proves the finite-pair trace law
for a normalized measurement tube, controls arbitrary legal quadratic
remainders, and uses a readout independent of that unknown remainder.
The factor $\sqrt{Nb}$ comes with both a block-state Bures upper and
an integer-block Bernoulli lower, not just a Fisher-information scaling.


\subsection{Memory-constrained testers and retained receiver systems}
Ohst--Zhang--Nguyen--Pl\'avala--Quintino \cite{OZNPQ89}
formulate discrimination with prescribed auxiliary dimensions as
optimization over constrained separable operators.  Their
Definition~6 specifies single-call memory-limited testers, while
Theorem~9 characterizes convex combinations of memoryless testers
by a convergent symmetric-extension hierarchy;
Sections~5.2--5.4 treat multi-call quantum and classical memory
and affine constraints on the factors of the lifted decomposition.
These are dimension and causal-transfer restrictions, not a bound on
the squared number of calls between fresh preparations.

Their Definition~23, equation~(64), describes two-call classically
adaptive testers as $T_i=\sum_zR_z\otimes S_{i|z}$, with single-call
testers $R$ and $S_{\cdot|z}$.  The first complete output is measured
and only the classical result chooses the second tester.  Their
Theorem~24 shows that this class and the parallel class are not
nested; both lie in the adaptive class.  Our inclusion and explicit
negative-partial-transpose witness in
Proposition~\ref{prop:memorycomparison89} answer the relevant
comparison without identifying these classical testers with all
unit-reset protocols.  Reset protocols may keep an unmeasured old
receiver register for a final joint measurement.  It is independent
of the complete next acquisition, but need not be separable from
all other receiver systems after receiver processing.

The constrained-separability formulation in their Section~6.2,
equations~(65)--(68), optimizes over the stated classically adaptive
set.  It is not a characterization of our larger retained-receiver
set: its cross-use separability requirement excludes the tester
\eqref{eq:memorywitness89}.  Conditional tensor independence at
preparation is not the same condition as separability of every
final tester effect across uses.  Consequently an upper bound for
that smaller optimization cannot be used as an upper for the reset
class merely by dropping a dimension cap.  Their general adaptive
optimization, when applicable to the same finite channel ensemble,
can of course bound a reset subclass; that inclusion does not
supply a sharp $Q$-dependent shrinking-tube law.  We do not claim
that no other lifted constrained-separability representation of
reset preparation could be developed.

Thus arbitrary receiver memory with zero coherent import has an
explicit operational relation to their complete-measurement
class, but it is not identified with any fixed auxiliary-dimension
slice of their hierarchy.  The contribution here is not a new
memory hierarchy or a new separability method.  It is the
support-coordinate local profile and hard-budget theorem, including
the unknown quadratic remainder and a matching lower in those
resources.  The Bell-reference witness alone establishes a model
boundary; Theorem~\ref{thm:operationalhierarchy90} now supplies a strict
exact score hierarchy for the specified finite ensemble, whereas the general memory-constrained discrimination
problem and independent priority assessment remain distinct.

\subsection{Channel Bures bounds and block testing}
Huang--Meyer--Nuradha--Wilde \cite[Theorem~17, Remark~18]{HMNW88}
restate Yuan--Fung's parallel bound \cite[equation~(23)]{YF76}.
For aligned canonical isometries put
$M_W=V_0^*(W\otimes I)V_1$, where $\|W\|\le1$,
$a_W=\|I-\operatorname{Re}M_W\|$ and $b_W=\|I-M_W\|$.
Their Theorems~17 and~19 give, respectively,
\begin{align*}
 d_B^2(\mathcal N_0^{\otimes n},\mathcal N_1^{\otimes n})
 &\le n\inf_W\{2a_W+(n-1)b_W^2\},\\
 \sup_{\text{adaptive controls}}d_B^2(\mathcal P_0^{(n)},\mathcal P_1^{(n)})
 &\le n\inf_W\{2a_W+(n-1)\sqrt{2a_W}\,b_W\}.
\end{align*}
Channel Bures distance takes the supremum over complete
probe--reference inputs.  For our chosen surrogate dilations,
$2a_W=\|W_s-W_0\|^2\le\alpha^2s^2$ and
$b_W\le\alpha s$; either expression yields
$d_B\le\alpha ns$.  Restoring the actual tube member costs
$nrs^2$ in diamond norm and hence $s\sqrt{rn}$ in Bures distance.
Thus the block upper specializes established isometric-extension
bounds; it is not a new fidelity principle.  The finite-tube
contribution is the uniform remainder, complete support cases and
matching width-dependent lower.  The reset theorem additionally
handles history-dependent allocation and receiver instruments by
conditional fidelity, rather than unconditional tensorization.

Ghosal--Halder--Patra--Sen \cite[Section~III, Theorem~1]{GHPS88}
use block-i.i.d. probes, with arbitrary within-block correlations
and independent identical repetition.  For a closed convex null
channel set $\mathsf R$ and a simple alternative $\mathcal N$,
their memory-assisted fixed-block Stein exponent is
\[
 \sup_{\sigma_{RA^{\otimes\ell}}}\inf_{\Phi\in\mathsf R}
 \frac1\ell D\!\left((\operatorname{id}_R\otimes\Phi^{\otimes\ell})(\sigma)
 \middle\|(\operatorname{id}_R\otimes\mathcal N^{\otimes\ell})(\sigma)\right).
\]
Their memory-free expression restricts $\sigma$ to the probes.
The block hierarchy is a direct antecedent of our resource taxonomy.
Our blocks may have unequal sizes and nonidentical states; the reset
class also has classical feedback and adaptive within-block control.
We consider finite unhalved trace separation in a shrinking fixed tube,
not a composite Stein exponent at fixed block size.  Neither theorem
is substituted for the other, and no first introduction of block
strategies is claimed.

\subsection{Semidefinite tangent geometry and nonemptiness}
The positive-semidefinite tangent condition $QHQ\succeq0$ and
second-order corrections are classical variational geometry;
Bonnans--Cominetti--Shapiro \cite{BCS87} study parabolic tangent
sets and semidefinite second-order regularity.  We use no new
primitive cone theorem.  Proposition~\ref{prop:nonempty87}
provides a common normalization, explicit coefficient and interval,
and a simultaneous operator-norm remainder bound for the measurement
tuple.  The sufficient bound can be large and is not a minimal
feasibility threshold.  Its exact certificate does not determine
optimal discrimination constants or synthesize quantum controls.

\subsection{Fisher frames and covariance embeddings}
Saini--Kiukas--Burgarth--Gilchrist \cite[Theorem~1]{SKBG84} compare
classical and quantum Fisher information for a varying full-rank
\emph{state} measured by one fixed informationally complete POVM,
using a state-dependent frame operator.  Here the varying object
is the measurement, the tangent has coupled normalization, and the
loss permits $N$ adaptive uses with a reference.  In the rank-one
case our support span equals the effect span, so informational
completeness excludes the linear orbit regime by
Corollary~\ref{cor:orbitexamples84}.  That intersection does not
identify the two quadratic forms or their resource models.

Mayer--Yun \cite[Theorem~3.6, Definition~3.7 and Remark~4.4]{MY84}
use a kernel embedding of operator-valued measures, a kernel-dependent
discrepancy, and a tomographic Gram operator whose nullspace records
unobserved state directions.  Our support span contains entire
matrix blocks $P_j\mathcal H_dP_j$, rather than one observable
per effect; this distinction disappears for rank-one effects but
not in general.  Theorem~\ref{thm:supportkernel84} concerns a
coupled measurement tangent, and \eqref{eq:frontupper84} concerns
adaptive finite-use loss.  Neither kernel embedding nor injectivity
alone supplies those statements.

\subsection{Known symmetry and tomography primitives}
The full primary text of Yoshida--Okigami--Posta--Grinko
\cite{YOPG84} distinguishes known-symmetry estimation from known-pair
discrimination.  Theorem~7 gives invariant-state copy complexity
$\Theta((D_{\rm st}-1+\log(1/\eta))/\epsilon^2)$, where
$D_{\rm st}=\sum_\lambda m_\lambda^2$; $D_{\rm st}=1$ is a
known singleton.  Corollary~9 gives parallel covariant-channel bounds
\[
 O((D'_G+K_G\log(1/\eta))/\epsilon_\diamond^2),\qquad
 O((C'_G+\log(1/\eta))/\epsilon_{\rm tr}^2),
\]
in one-use halved diamond distance and normalized-Choi trace distance,
respectively, for $0<\epsilon\le1$, $0<\eta\le1/2$.
Here $D_G=\sum_\lambda m_\lambda M_\lambda$ counts Choi commutant
parameters before trace preservation,
$K_G=\sum_\lambda m_\lambda$,
$D'_G=D_G+K_G\log(4L)$, and
$C'_G=d_{\mathcal I}^{-1}(\sum_\lambda m_\lambda
\sqrt{d_\lambda M_\lambda})^2$; $L$ counts occupied blocks.
Theorem~13 does not give a general matching inverse-square diamond
lower: its bound is $\Omega_\xi(D_G/\epsilon_\diamond^{\xi/2})$
under its affine-dimension hypothesis.

For permutation covariance on $\ell$ systems of fixed local dimension
$q$, Theorems~12 and~15 match the Choi order
$\ell^{q^4-q^2}\epsilon_{\rm tr}^{-2}$ at constant confidence.
Theorem~18 approximates the Hayashi measurement in
$O((s+q)\operatorname{polylog}(s,q,\tau^{-1},\zeta^{-1}))$ gates,
with coupled output error at most $\tau$ except with probability
$\zeta$.  This is not a diamond-error certificate for our recovery.
Known-symmetry estimation, its parameter dimensions, and that
specialized gate construction do not themselves give our support
kernel or finite-pair curve theorem.  Nor do our theorems claim
faster learning on those symmetry classes.

For ordinary tomography the main-text substitutions remain explicit.
Mele--Bittel \cite[Theorem~III.3]{MB67} give the sufficient
constant-confidence order $O(dkr\epsilon^{-2})$ in halved diamond
loss, with $r\le dk$ for measurement channels.
Zambrano--Ramos-Calderer--Kueng \cite[Theorem~2]{ZRK76} give the
sufficient single-copy worst-input outcome-TV bound
\[
 m\ge\frac{8[d^2+\epsilon(d^2+1)/6]}{\epsilon^2}
          \log\frac{2^{k+1}9^{2d}}{\eta}.
\]
Either loss bounds the component norm error by $\epsilon$ after
output dephasing where needed.  A legal estimator need not be
balanced: apply Theorem~\ref{thm:affinerepair83} before using the
interior metric.  The resulting future loss is at most
$8k\sqrt N\epsilon$.  With $\epsilon=\delta/(8k\sqrt N)$ the
respective sufficient orders are
\[
 O(d^2k^4N\delta^{-2}),\qquad
 O\!\left(k^2N\delta^{-2}d^2[d+k+\log(1/\eta)]\right).
\]
These are upper-bound substitutions, not optimality claims about
the cited estimators.  Our component procedure credits the existing
binary operator-norm primitive and uses a different confidence
allocation.  The full calculations and confidence restrictions
are retained in the supplement, Section~\ref{sec:priority83}.

\subsection{What remains beyond the theorems}
The orbit theorem and Theorem~\ref{thm:curveclassification85} provide
matching scales on each fixed two-sided smooth curve with nonzero
first derivative, including second-order rank openings.  Their
constants need not be uniform across curves.  Arbitrary-pair
midpoint equivalence and full-body multi-outcome entropy are not
consequences of these local statements.  Growing-$k$ minimax sharpness,
polynomial entropy-optimal dictionaries, and general readout synthesis
remain separate mathematical questions.  Independent expert priority
assessment is not inferred from the absence of an identical statement
in this targeted comparison.

\subsection{Tangent neighborhoods and independent acquisition}
Theorem~\ref{thm:tubedichotomy86} fixes $E,H$ and an explicit
second-order remainder allowance, rather than assuming a particular
smooth family of Kraus operators.  It treats the full one-sided
positive tangent cone.  A nonzero missing-support compression has a
classical impossible-event witness; only within the cone's lineality
space does the support-generator test select the coherent branch.
The latter still uses the credited correction mechanism.
Theorem~\ref{thm:producttrichotomy86} separates these mechanisms by
an explicit product probe--reference restriction.  Its upper uses
standard root-fidelity inequalities and tensor multiplicativity
\cite[Chapter~3]{Watrous65}, not a new fidelity principle or a bound
on all nonadaptive entangled strategies.  The intervening parallel
class is addressed instead by the new proof in
Section~\ref{sec:width87}.  Its conclusion matches the adaptive
local order and does not promote the product upper to that class.

The higher-order statement requires the finite remainder
\eqref{eq:powerjet86}.  Proposition~\ref{prop:mixedrates86} instead
proves a two-scale law when an intermediate support-opening term is
present.  These results address one-sided openings and selected
higher-order approaches without identifying a first jet with an
arbitrary-pair global geometry.  No general boundary entropy or
unrestricted growing-outcome learning law is inferred.

\subsection{Autonomous coherent-memory hierarchies}
Zonnios--Binder \cite[Proposition~2]{ZB90} prove monotonicity of their
memory-parametrized process distinguishability and its upper bound by
the strategy norm.  Their Theorem~1 compiles any finite-time tester into
repeated use of one instrument with an internal counter and sufficient
coherent memory; Corollary~1 gives finite-time completeness.  Their
unknown object is a multi-time process, and their resource is the coherent
ancilla dimension of an autonomous repeated-instrument machine.
Our reset class instead allows public time-dependent history rules and
unrestricted receiver storage, but forbids coherent import across a
specified fresh-preparation cut.  Its hard square cost $Q$ is not an
ancilla dimension.  Neither their completeness theorem nor our fixed-tube
profile law supplies the other's conclusion.  The exact finite-ensemble
hierarchy of Theorem~\ref{thm:operationalhierarchy90} likewise compares
three acquisition/receiver classes, not every dimension slice of their
hierarchy.  Their work is a direct antecedent for systematic coherent-memory
interpolation; no first such framework is claimed here.

\subsection{Classical outputs and accessible post-measurement states}
Eid--Quintino \cite{EQ90} consider the quantum state returned by the
associated L\"uders instrument in addition to the classical label.
Our unknown channel returns only that label.  The references used in
Theorem~\ref{thm:operationalhierarchy90} were supplied by the experimenter;
they are not residual quantum outputs of the device.  Thus an instrument
discrimination advantage cannot be imported into the present interface.

\subsection{What the exact hierarchy adds}
The strict tester inclusion in Proposition~\ref{prop:memorycomparison89}
alone is not an operational score theorem.  Theorem~\ref{thm:operationalhierarchy90}
provides one finite classical-output ensemble, exact upper bounds for all
three classes, and matching realizations in every dimension and along a
full-rank deformation.  Ohst et al.\ \cite[Theorem~24]{OZNPQ89} already
separate other strategy classes by channel-discrimination examples; the
present assertion does not claim a first memory advantage.  Its middle
optimization permits arbitrary unmeasured receiver storage and all
unit-reset feedback.  Lemma~\ref{lem:swapfilter90} is the step that prevents
a separable-tester upper from being incorrectly assigned to this larger
class.  These exact ensemble values and the local allocation orders
answer different questions; neither is substituted for the other.
